\documentclass[11pt]{article}
\usepackage[a4paper,margin=27mm]{geometry}
\usepackage{amsmath,amssymb,amsthm,xcolor,booktabs,array}
\newcommand{\PROVED}{\textcolor{green!40!black}{\textbf{PROVED}}}
\newcommand{\OPEN}{\textcolor{red!70!black}{\textbf{OPEN}}}
\newcommand{\AUDIT}{\textcolor{orange!70!black}{\textbf{AUDIT}}}
\newtheorem{theorem}{Theorem}
\newtheorem{proposition}{Proposition}
\title{Exact Tate renormalization of the Poisson operator (v106)}
\date{13 August 2026}

\begin{document}
\maketitle

\section{The rank-two subtraction}

Let $f$ be an even Schwartz function on $\mathbb R$, let
\[
 Zf(x)=\sum_{n\geq1}f(nx),\qquad
 (\mathscr Jg)(x)=x^{-1}g(x^{-1}),
\]
and use the self-dual additive Fourier transform.  Poisson summation gives
the exact identity
\begin{equation}
 f(0)+2Zf(x)=x^{-1}\widehat f(0)
                 +2x^{-1}Z\widehat f(x^{-1}).                 \tag{1}
\end{equation}
Define the polar operator and the renormalized operator by
\begin{equation}
 \Pi f(x)=\frac{\widehat f(0)}{2x}-\frac{f(0)}2,
 \qquad \mathcal Rf=Zf-\Pi f.                                \tag{2}
\end{equation}

\begin{theorem}[Exact Gamma--Fourier--Tate identity --- \PROVED]
For every even Schwartz function,
\begin{equation}
                 \boxed{\mathcal R=\mathscr JZ\mathcal F}.   \tag{3}
\end{equation}
Moreover $\Pi$ and $\mathcal R$ commute with every arithmetic dilation
$U_nf(x)=f(nx)$.
\end{theorem}

\begin{proof}
Solving (1) for $Zf$ gives (3).  Since
\[
 (U_nf)(0)=f(0),\qquad \widehat{U_nf}(0)=n^{-1}\widehat f(0),
\]
we have
\[
 \Pi(U_nf)(x)=\frac{\widehat f(0)}{2nx}-\frac{f(0)}2
             =U_n(\Pi f)(x).
\]
The operator $Z$ commutes with $U_n$, first on finite sums and then in
Meyer's multiplier topology, so $\mathcal R=Z-\Pi$ also commutes with
$U_n$.
\end{proof}

The formula is the desired subtraction at the level of the pole extension:
the two coefficients are exactly the two Tate evaluations $f(0)$ and
$\widehat f(0)$.  It is important, however, to determine what it changes
on the odd object rather than infer this from the rank of $\Pi$.

\section{Effect on the odd quotient}

Put
\begin{equation}
 \mathcal H_\cap=\{f\in\mathcal H_+:f(0)=\widehat f(0)=0\}.   \tag{4}
\end{equation}

\begin{theorem}[Tate subtraction leaves the odd image unchanged ---
\PROVED]
On $\mathcal H_\cap$ one has
\begin{equation}
                       \mathcal Rf=Zf.                        \tag{5}
\end{equation}
Consequently
\begin{equation}
 \mathcal R\mathcal H_\cap=Z\mathcal H_\cap,
 \qquad
 \mathcal H_-/\mathcal R\mathcal H_\cap
       =\mathcal H_-/Z\mathcal H_\cap                       \tag{6}
\end{equation}
as topological scaling representations.  In particular, the nuclear
character and every finite Loewy subquotient are exactly the same as in
row (c).
\end{theorem}

\begin{proof}
Both coefficients in (2) vanish on (4), proving (5), hence equality of
the two closed invariant images.  The identity map on $\mathcal H_-$ then
induces the identity of the quotients in (6), including their topology,
scaling action, nuclear trace and Loewy filtration.
\end{proof}

Thus the rank-two renormalization is not a new metric on the odd quotient.
It is an exact splitting calculation for the even polar extension.  Any
Hilbert norm on the quotient obtained after (2) still has to prove the
same centered Witt estimate
\begin{equation}
 \bigl\|\sqrt n\,U_n\bigr\|_{\operatorname{gr}_L M}\leq1     \tag{7}
\end{equation}
for every finite odd Loewy factor $M$.

\section{Why no endpoint norm is hidden in (3)}

For a general even Schwartz function, $\mathcal Rf$ need not lie in an
unweighted $L^p$ space.  Indeed, $Zf(x)\to0$ as $x\to\infty$, whereas
from (2)
\begin{equation}
 \mathcal Rf(x)=Zf(x)-\frac{\widehat f(0)}{2x}+\frac{f(0)}2
       \longrightarrow\frac{f(0)}2.                          \tag{8}
\end{equation}
The apparent residual constant is the opposite Tate chart in (3), not an
$L^p$ error.  Requiring both Tate evaluations to vanish removes it and
returns exactly $Z\mathcal H_\cap$, by (5).  Therefore interpolation of
$\mathcal R$ on an ambient unweighted $L^p$ cannot furnish the missing odd
Hilbert completion.

\begin{center}
\begin{tabular}{>{\raggedright\arraybackslash}p{.58\linewidth}
                >{\centering\arraybackslash}p{.15\linewidth}
                >{\raggedright\arraybackslash}p{.17\linewidth}}
\toprule
Statement & Status & Location \\
\midrule
Exact Poisson--Tate subtraction & \PROVED & (1)--(3) \\
Covariance with all Witt dilations & \PROVED & Theorem 1 \\
Equality of renormalized and original odd quotients & \PROVED & (5)--(6) \\
Ambient unweighted endpoint bound & false in general & (8) \\
Contractivity (7), condition $G$, row (d) & \OPEN & unchanged on the odd quotient \\
\bottomrule
\end{tabular}
\end{center}

\end{document}
